\subsubsection{Expectation Value}

In the previous sections, we have seen how to compute the probability of obtaining a certain measurement outcome when measuring an observable on a quantum state. However, in many cases, we are \hl{interested in the average value of an observable over many measurements}, rather than the probability of specific outcomes. This average value is known as the \textbf{expectation value} of the observable.

\highspace
\begin{definitionbox}[: Expectation Value]
    Let $O$ be an observable represented by a Hermitian operator, and let \ket{\psi} be a quantum state. The quantity:
    \begin{equation}\label{eq:expectation-value}
        \langle O \rangle = \bra{\psi} O \ket{\psi} = \sum_k \lambda_k p_k
    \end{equation}
    Is called the \definition{Expectation Value} of the observable $O$ in the state \ket{\psi}.
\end{definitionbox}

\highspace
\textcolor{Green3}{\faIcon{book} \textbf{Physical Meaning.}} The expectation value $\langle O \rangle$ can be interpreted as the \hl{average result obtained by measuring the observable $O$ many times on identical copies of the state $\ket{\psi}$.}
\begin{enumerate}
    \item Prepare many identical copies of the state \ket{\psi}.
    \item Measure the observable $O$ on each copy.
    \item Record the outcomes.
    \item Compute the arithmetic mean of the recorded outcomes.
\end{enumerate}
As the number of experiments increases, the average converges to the expectation value $\bra{\psi} O \ket{\psi}$.

\highspace
\begin{flushleft}
    \textcolor{Red2}{\faIcon{exclamation-triangle} \textbf{Important Clarification}}
\end{flushleft}
The expectation value is \textbf{not necessarily} one of the possible measurement outcomes. It is the \textbf{average (more precisely, the probability-weighted average)} of many outcomes, and can take values that are not eigenvalues of the observable. For example, if possible outcomes are $+1$ with probability 50\% and $-1$ with probability 50\%, then the expectation value is:
\begin{equation*}
    \left(+1\right) \cdot \frac{1}{2} + \left(-1\right) \cdot \frac{1}{2} = 0
\end{equation*}
Although the actual measurement never returns $0$, the expectation value is $0$ because it represents the average of many measurements. This means that half of the measurements return $+1$ and the other half return $-1$, resulting in an average of $0$.

\highspace
The expectation value tells us the ``center of mass'' of the measurement distribution. \hl{It summarizes the statistical behavior of the measurement in a single number.}

\highspace
\textcolor{Green3}{\faIcon{question-circle} \textbf{Why are we interested in the average value of an observable rather than in the probability of specific outcomes?}} Because in many real situations we do not \textbf{care about} the result of a single measurement, but about the \textbf{overall behavior of the system}. Also, it proves a stable, experimentally meaningful quantity that summarizes the statistical behavior of an observable over many measurements, even when individual outcomes are random.

\highspace
\begin{flushleft}
    \textcolor{Green3}{\faIcon{tools} \textbf{How can we compute the expectation value?}}
\end{flushleft}
To compute the expectation value $\langle O \rangle$ of an observable $O$ in a state \ket{\psi}, we start from the Spectral Theorem (\autopageref{eq:spectral-decomposition}):
\begin{equation*}
    O = \sum_k \lambda_k \proj{k}
\end{equation*}
Substitute this into the expectation value formula (\autopageref{eq:expectation-value}):
\begin{align*}
    \bra{\psi} O \ket{\psi} &= \bra{\psi} \left( \sum_k \lambda_k \proj{k} \right) \ket{\psi} \\
    &\downarrow
    \text{ use linearity of the inner product} \\
    &=
    \sum_k \lambda_k \braket{\psi}{k} \braket{k}{\psi}
\end{align*}
Using the property of inner products (\autopageref{sec:magnitude}):
\begin{equation*}
    \braket{\psi}{k} = \overline{\braket{k}{\psi}} \quad \implies \quad \overline{\braket{k}{\psi}} \braket{k}{\psi} = \left|\braket{k}{\psi}\right|^{2}
\end{equation*}
Here, the Born rule (\autopageref{eq:born-rule}) states that $\left|\braket{k}{\psi}\right|^{2}$ is the probability of measuring the outcome $\lambda_k$ when measuring $O$ on the state \ket{\psi}. Therefore, we can rewrite the expectation value as:
\begin{equation}\label{eq:expectation-value-probabilities}
    p_k = \left|\braket{k}{\psi}\right|^{2} \quad \implies \quad \bra{\psi} O \ket{\psi} = \sum_k \lambda_k p_k
\end{equation}
This formula shows that the expectation value is the \textbf{weighted average} of the possible outcomes $\lambda_k$, where each outcome is weighted by its probability $p_k$. In other words, the \hl{expectation value is the sum of all possible outcomes multiplied by their respective probabilities}, which is consistent with our interpretation of it as an average over many measurements.